\section*{Abstract}

Scientific workflows increasingly produce structured \gls{json} data that is easy to exchange but difficult to interpret semantically across systems. While JSON Schema supports structural validation, it does not provide native semantics for Linked Data integration. This thesis addresses that gap by extending MetaConfigurator, a schema-driven web application for structured data authoring, with Semantic Web capabilities.

The core contribution is an integrated \gls{rdf} Authoring Panel that supports end-to-end semantic data workflows in a single User Interface. It enables researchers to connect existing traditional hierarchical or tabular research data to the Semantic Web by using AI-assisted \gls{rml} to transform these data into \gls{rdf}. After transformation, the same panel supports refinement and modification of triples, \gls{sparql}-based exploration, interactive Knowledge Graph visualization, and export to common \gls{rdf} serializations for import into existing knowledge stores.

The result is a unified tool that bridges traditional non-linked research data and Semantic Web infrastructures. The thesis demonstrates this tool on a real-world \gls{mof} synthesis application, showing how tabular research data can be transformed to \gls{rdf}, iteratively improved, explored, visualized, and finally integrated into a target knowledge store.

\newpage

\section*{Kurzfassung}

Wissenschaftliche Workflows erzeugen zunehmend strukturierte \gls{json}-Daten, die leicht austauschbar sind, jedoch systemübergreifend schwer semantisch zu interpretieren sind. Während JSON Schema die strukturelle Validierung unterstützt, bietet es keine nativen Semantiken für die Integration von Linked Data. Diese Arbeit schließt diese Lücke, indem sie MetaConfigurator, eine schema-gesteuerte Webanwendung zur Erstellung strukturierter Daten, um Semantic-Web-Funktionalitäten erweitert.

Der zentrale Beitrag ist ein integriertes \gls{rdf}-Authoring-Panel, das End-to-End-Workflows für semantische Daten in einer einzigen Benutzeroberfläche unterstützt. Es ermöglicht Forschenden, bestehende traditionelle hierarchische oder tabellarische Forschungsdaten mit dem Semantic Web zu verknüpfen, indem KI-gestütztes \gls{rml} verwendet wird, um diese Daten in \gls{rdf} zu transformieren. Nach der Transformation unterstützt dasselbe Panel die Verfeinerung und Modifikation von Tripeln, die \gls{sparql}-basierte Exploration, die interaktive Visualisierung von Wissensgraphen sowie den Export in gängige \gls{rdf}-Serialisierungen zur Integration in bestehende Wissensspeicher.

Das Ergebnis ist ein einheitliches Werkzeug, das traditionelle, nicht verknüpfte Forschungsdaten mit Semantic-Web-Infrastrukturen verbindet. Die Arbeit demonstriert dieses Werkzeug anhand einer realen Anwendung zur \gls{mof}-Synthese und zeigt, wie tabellarische Forschungsdaten in \gls{rdf} transformiert, iterativ verbessert, exploriert, visualisiert und schließlich in einen Ziel-Wissensspeicher integriert werden können.

